
\documentclass[11pt]{article}
%%%%%%%%%%%%%%%%%%%%%%%%%%%%%%%%%%%%%%%%%%%%%%%%%%%%%%%%%%%%%%%%%%%%%%%%%%%%%%%%%%%%%%%%%%%%%%%%%%%%%%%%%%%%%%%%%%%%%%%%%%%%%%%%%%%%%%%%%%%%%%%%%%%%%%%%%%%%%%%%%%%%%%%%%%%%%%%%%%%%%%%%%%%%%%%%%%%%%%%%%%%%%%%%%%%%%%%%%%%%%%%%%%%%%%%%%%%%%%%%%%%%%%%%%%%%
\usepackage{makeidx}
\usepackage{amsfonts}
\usepackage{amssymb}
\usepackage{amsmath}

\setcounter{MaxMatrixCols}{10}
%TCIDATA{OutputFilter=LATEX.DLL}
%TCIDATA{Version=5.50.0.2953}
%TCIDATA{<META NAME="SaveForMode" CONTENT="1">}
%TCIDATA{BibliographyScheme=Manual}
%TCIDATA{Created=Monday, September 29, 2008 22:43:53}
%TCIDATA{LastRevised=Sunday, September 15, 2024 14:54:10}
%TCIDATA{<META NAME="GraphicsSave" CONTENT="32">}
%TCIDATA{<META NAME="DocumentShell" CONTENT="Standard LaTeX\Blank - Standard LaTeX Article">}
%TCIDATA{CSTFile=40 LaTeX article.cst}

\newtheorem{theorem}{Theorem}
\newtheorem{acknowledgement}[theorem]{Acknowledgement}
\newtheorem{algorithm}[theorem]{Algorithm}
\newtheorem{axiom}[theorem]{Axiom}
\newtheorem{case}[theorem]{Case}
\newtheorem{claim}[theorem]{Claim}
\newtheorem{conclusion}[theorem]{Conclusion}
\newtheorem{condition}[theorem]{Condition}
\newtheorem{conjecture}[theorem]{Conjecture}
\newtheorem{corollary}[theorem]{Corollary}
\newtheorem{criterion}[theorem]{Criterion}
\newtheorem{definition}[theorem]{Definition}
\newtheorem{example}[theorem]{Example}
\newtheorem{exercise}[theorem]{Exercise}
\newtheorem{lemma}[theorem]{Lemma}
\newtheorem{notation}[theorem]{Notation}
\newtheorem{problem}[theorem]{Problem}
\newtheorem{proposition}[theorem]{Proposition}
\newtheorem{remark}[theorem]{Remark}
\newtheorem{solution}[theorem]{Solution}
\newtheorem{summary}[theorem]{Summary}
\newenvironment{proof}[1][Proof]{\noindent\textbf{#1.} }{\ \rule{0.5em}{0.5em}}
\newenvironment{question}[1][Question]{}{}
\newenvironment{answer}[1][Answer]{}{}
\newenvironment{draftnote}[1][Draft note]{}{}
\input{tcilatex}
\begin{document}


\begin{draftnote}
\pprintfilename
\end{draftnote}

\begin{exercise}
\cite{ljungqvist_sargent:2020} Exercise 7.11. \textbf{Pissarides' analysis
of taxation and variable search intensity}. At the beginning of each period,
an unemployed worker chooses search effort $e$, and receives a wage offer
with probability $\pi \left( e\right) $ and no offer with probability $1-\pi
\left( e\right) $. The search effort---a measure of search
intensity---incurs a period monetary cost $e$. The wage offer $w$ is drawn
from distribution $F$, and the worker chooses to accept or reject it. If the
offer is rejected or no offer is received, the worker receives unemployment
benefit $z>0$, and starts as an searching unemployed worker at the beginning
of the next period. If the wage is accepted, the worker receives wage $w$
minus income tax $T\left( w\right) $, and continues as an employed worker
without searching. At the beginning of each period, a worker who was
employed in the previous period can be fired with probability $\alpha $, $%
0<\alpha <1$, in which case the worker receives unemployment benefits $z$ in
the given period and becomes a searching unemployed worker at the beginning
of the next period.

The worker chooses a strategy, that is, whether to accept an offer or not
and how much effort to put into search when unemployed, to maximize%
\begin{equation*}
E\sum_{t=0}^{\infty }\beta ^{t}y_{t},\qquad 0<\beta <1
\end{equation*}%
where $y_{t}$ is period income (either after-tax wage or unemployment
benefit) net of search costs. We assume that $w-T\left( w\right) $ is
increasing in $w$, $\pi \left( 0\right) =0$---that is, the worker gets no
offers without search effort---and that $\pi ^{\prime }\left( e\right) >0$, $%
\pi ^{\prime \prime }\left( e\right) >0$.

\begin{enumerate}
\item[a)] Analyze the worker's problem. Establish that the optimal strategy
is to choose a reservation wage. Display the condition that describes the
optimal choice of $e$, and show that the reservation wage is independent of $%
e$.

\item[b)] Assume that $T\left( w\right) =\tau \left( w-a\right) $ where $%
0<\tau <1$ and $a>0$, i.e., the tax schedule consists of a flat tax rate $%
\tau $ and a lump-sum rebate $\tau a$. Show that an increase in $a$
decreases the reservation wage and increases the level of effort, increasing
the probability of accepting employment.

\item[c)] Show under what conditions a change in $\tau $ has the opposite
effect.
\end{enumerate}

\begin{solution}
For \textbf{question a)}, the value of a searching worker, before the cost
of search has been paid and offer revealed, is%
\begin{equation}
Q=\max_{e}-e+\pi \left( e\right) \int_{0}^{B}V\left( w^{\prime }\right)
dF\left( w^{\prime }\right) +\left( 1-\pi \left( e\right) \right) \left(
z+\beta Q\right)   \label{eq:Pissarides_Q}
\end{equation}%
and the value of an unemployed worker with offer $w$ at hand, after the cost
has been paid, is%
\begin{equation}
V\left( w\right) =\max_{\left\{ \text{accept, reject}\right\} }\left\{
V^{a}\left( w\right) ,z+\beta Q\right\} .  \label{eq:Pissarides_V}
\end{equation}%
The value of an accepted offer is%
\begin{equation*}
V^{a}\left( w\right) =w-T\left( w\right) +\beta \left[ \alpha \left( z+\beta
Q\right) +\left( 1-\alpha \right) V^{a}\left( w\right) \right] \text{.}
\end{equation*}%
This is identical to the value of an employed worker with wage $w$ who has
not been fired at the beginning of the period. Hence%
\begin{equation*}
V^{a}\left( w\right) =\frac{1}{1-\beta \left( 1-\alpha \right) }\left[
w-T\left( w\right) +\beta \alpha \left( z+\beta Q\right) \right] 
\end{equation*}%
which is strictly increasing in $w$ because $w-T\left( w\right) $ is
strictly increasing in $w$. Here we assumed that a fired worker stays
unemployed for one period without searching and then rejoins search in the
subsequent period. Equation\ (\ref{eq:Pissarides_V}) for $V\left( w\right) $
implies that the optimal policy is a reservation wage---the value of
accepting is strictly increasing in $w$, while the value of rejecting is
independent of $w$.\ Equation (\ref{eq:Pissarides_V}) implies that the
reservation wage $\bar{w}$ is given implicitly by%
\begin{equation*}
\frac{1}{1-\beta \left( 1-\alpha \right) }\left[ \bar{w}-T\left( \bar{w}%
\right) +\beta \alpha \left( z+\beta Q\right) \right] =1+z+\beta Q,
\end{equation*}%
and hence%
\begin{equation}
\frac{\bar{w}-T\left( \bar{w}\right) }{1-\beta }=z+\beta Q=V^{a}\left( \bar{w%
}\right) =V\left( \bar{w}\right) .  \label{eq:Pissarides_wbar}
\end{equation}%
Thus $\bar{w}$ is independent of the currrent period choice $e$ (even though
the choice of future $e$ implicitly shows up in $Q$). Current-period choice
of $e$ is irrelevant because that effort decision is already sunk when
making the acceptance choice.

The optimal choice of $e$ is given by the first-order condition to (\ref%
{eq:Pissarides_Q}):%
\begin{equation*}
1=\pi ^{\prime }\left( e^{\ast }\right) \left[ \int_{0}^{B}V\left( w^{\prime
}\right) dF\left( w^{\prime }\right) -\left( z+\beta Q\right) \right] 
\end{equation*}%
Further notice that%
\begin{eqnarray}
\int_{0}^{B}V\left( w^{\prime }\right) dF\left( w^{\prime }\right) 
&=&\int_{0}^{\bar{w}}\frac{\bar{w}-T\left( \bar{w}\right) }{1-\beta }%
dF\left( w^{\prime }\right) +\int_{\bar{w}}^{B}V^{a}\left( w^{\prime
}\right) dF\left( w^{\prime }\right)   \label{eq:Pissarides_EV} \\
&=&\frac{\bar{w}-T\left( \bar{w}\right) }{1-\beta }+\frac{1}{1-\beta \left(
1-\alpha \right) }\int_{\bar{w}}^{B}\left[ w^{\prime }-T\left( w^{\prime
}\right) -\left( \bar{w}-T\left( \bar{w}\right) \right) \right] dF\left(
w^{\prime }\right) ,  \notag
\end{eqnarray}%
and hence the optimality condition for the choice of effort reads%
\begin{equation}
1=\frac{\pi ^{\prime }\left( e^{\ast }\right) }{1-\beta \left( 1-\alpha
\right) }\int_{\bar{w}}^{B}\left[ w^{\prime }-T\left( w^{\prime }\right)
-\left( \bar{w}-T\left( \bar{w}\right) \right) \right] dF\left( w^{\prime
}\right)   \label{eq:Pissarides_opt1}
\end{equation}%
which pins down a relationship between optimal effort and reservation wage.
In particular, a higher reservation wage decreases the integral, so $\pi
^{\prime }\left( e^{\ast }\right) $ has to increase, so $e^{\ast }$ has to
decrease (since $\pi ^{\prime \prime }\left( e\right) <0$). More formally,
using the implicit function theorem on (\ref{eq:Pissarides_opt1}) we obtain%
\begin{eqnarray}
\frac{de^{\ast }}{d\bar{w}} &=&-\frac{\partial f/\partial \bar{w}}{\partial
f/\partial e^{\ast }}=\frac{\pi ^{\prime }\left( e^{\ast }\right) }{\pi
^{\prime \prime }\left( e^{\ast }\right) }\frac{\left( 1-T^{\prime }\left( 
\bar{w}\right) \right) \left( 1-F\left( \bar{w}\right) \right) }{\int_{\bar{w%
}}^{B}\left[ w^{\prime }-T\left( w^{\prime }\right) -\left( \bar{w}-T\left( 
\bar{w}\right) \right) \right] dF\left( w^{\prime }\right) }  \notag \\
&=&\frac{\left( \pi ^{\prime }\left( e^{\ast }\right) \right) ^{2}}{\pi
^{\prime \prime }\left( e^{\ast }\right) \left( 1-\beta \left( 1-\alpha
\right) \right) }\left( 1-T^{\prime }\left( \bar{w}\right) \right) \left(
1-F\left( \bar{w}\right) \right) <0.  \label{eq:Pissarides_dedw}
\end{eqnarray}

Finally, from (\ref{eq:Pissarides_Q}), given optimal $e^{\ast }$, we can
express $Q$ as%
\begin{equation*}
Q=-e^{\ast }+\pi \left( e^{\ast }\right) \int_{0}^{B}V\left( w^{\prime
}\right) dF\left( w^{\prime }\right) +\left( 1-\pi \left( e^{\ast }\right)
\right) \left( z+\beta Q\right) 
\end{equation*}%
and hence%
\begin{equation*}
Q=\left[ 1-\beta \left( 1-\pi \left( e^{\ast }\right) \right) \right] ^{-1}%
\left[ -e^{\ast }+\pi \left( e^{\ast }\right) \int_{0}^{B}V\left( w^{\prime
}\right) dF\left( w^{\prime }\right) +\left( 1-\pi \left( e^{\ast }\right)
\right) z\right] 
\end{equation*}%
where we can substitute out $\int_{0}^{B}V\left( w^{\prime }\right) dF\left(
w^{\prime }\right) $ using (\ref{eq:Pissarides_EV}). Using this expression
for $Q$ in the indifference condition (\ref{eq:Pissarides_wbar}) yields%
\begin{eqnarray*}
\frac{\bar{w}-T\left( \bar{w}\right) }{1-\beta } &=&z+\beta Q \\
&=&\frac{z-\beta e^{\ast }}{1-\beta \left( 1-\pi \left( e^{\ast }\right)
\right) }+\frac{\beta \pi \left( e^{\ast }\right) }{1-\beta \left( 1-\pi
\left( e^{\ast }\right) \right) }\int_{0}^{B}V\left( w^{\prime }\right)
dF\left( w^{\prime }\right) 
\end{eqnarray*}%
and, substituting out $\int_{0}^{B}V\left( w^{\prime }\right) dF\left(
w^{\prime }\right) $, we obtain,%
\begin{equation}
\bar{w}-T\left( \bar{w}\right) -\left( z-\beta e^{\ast }\right) =\frac{\beta
\pi \left( e^{\ast }\right) }{1-\beta \left( 1-\alpha \right) }\int_{\bar{w}%
}^{B}\left[ w^{\prime }-T\left( w^{\prime }\right) -\left( \bar{w}-T\left( 
\bar{w}\right) \right) \right] dF\left( w^{\prime }\right) 
\label{eq:Pissarides_opt2}
\end{equation}%
To summarize, equations (\ref{eq:Pissarides_opt1}) and (\ref%
{eq:Pissarides_opt2}) yield two optimality conditions for unknowns $e^{\ast }
$ and $\bar{w}$.

Equation (\ref{eq:Pissarides_opt1}) compares the marginal cost of extra
effort on the left-hand side (which is equal to one, see equation (\ref%
{eq:Pissarides_Q})) with the marginal benefit on the right-hand side, which
is the present value of wage flows from future offers above reservation
wage, times marginal increase in probability $\pi ^{\prime }\left( e^{\ast
}\right) $ that these offers will arrive with extra unit of more effort.

Equation (\ref{eq:Pissarides_opt2}) is the familar optimality condition for
the choice of reservation wage. The left-hand side is the cost of searching
one more time when the wage offer at hand is equal to $\bar{w}$, which
consists of the net wages $\bar{w}-T\left( \bar{w}\right) $ above the flow
value of unemployment benefits net of the cost of search next period again, $%
z-\beta e^{\ast }$. The right-hand side is the value of search one more
time, which again consists of the net value of wage flows from offers above
the reservation wage. \ Finally notice that we can substitute out the
integral on the right-hand side of (\ref{eq:Pissarides_opt2}) using (\ref%
{eq:Pissarides_opt1}) to obtain%
\begin{equation*}
\bar{w}-T\left( \bar{w}\right) -\left( z-\beta e^{\ast }\left( \bar{w}%
\right) \right) =\frac{\beta \pi \left( e^{\ast }\left( \bar{w}\right)
\right) }{\pi ^{\prime }\left( e^{\ast }\left( \bar{w}\right) \right) }
\end{equation*}%
where we now formally write that the optimal choice of $e^{\ast }$ is a
function of $\bar{w}$, per (\ref{eq:Pissarides_opt1}). To deduce the
properties of the reservation wage, we can study the intersect(s) of the
functions%
\begin{eqnarray*}
h_{L}\left( w\right)  &=&w-T\left( w\right) -\left( z-\beta e^{\ast }\left(
w\right) \right)  \\
h_{R}\left( w\right)  &=&\frac{\beta \pi \left( e^{\ast }\left( w\right)
\right) }{\pi ^{\prime }\left( e^{\ast }\left( w\right) \right) }.
\end{eqnarray*}%
Differentiating both functions with respect to $w$, and using the expression
(\ref{eq:Pissarides_dedw}), we get%
\begin{eqnarray*}
\frac{d}{dw}h_{L}\left( w\right)  &=&\left( 1-T^{\prime }\left( w\right)
\right) \left[ 1+\beta \frac{\left( \pi ^{\prime }\left( e^{\ast }\left(
w\right) \right) \right) ^{2}}{\pi ^{\prime \prime }\left( e^{\ast }\left(
w\right) \right) }\frac{1-F\left( w\right) }{1-\beta \left( 1-\alpha \right) 
}\right]  \\
\frac{d}{dw}h_{R}\left( w\right)  &=&\left( 1-T^{\prime }\left( w\right)
\right) \left[ \beta \frac{\left( \pi ^{\prime }\left( e^{\ast }\left(
w\right) \right) \right) ^{2}}{\pi ^{\prime \prime }\left( e^{\ast }\left(
w\right) \right) }\frac{1-F\left( w\right) }{1-\beta \left( 1-\alpha \right) 
}-\beta \left( e^{\ast }\left( w\right) \right) \frac{1-F\left( w\right) }{%
1-\beta \left( 1-\alpha \right) }\right] 
\end{eqnarray*}%
where \thinspace $T^{\prime }\left( w\right) $ is the derivative of $T\left(
w\right) $.\ Comparing the two derivatives, we see that $\frac{d}{dw}%
h_{R}\left( w\right) <0$, and so the right-hand side curve is always
decreasing. On the other hand $\frac{d}{dw}h_{L}\left( w\right) $ could be
positive or negative but it is certainly true that%
\begin{equation*}
\frac{d}{dw}h_{L}\left( w\right) >\frac{d}{dw}h_{R}\left( w\right) .
\end{equation*}%
Therefore, $h_{L}\left( w\right) $ may not be increasing, but certainly
intersects $h_{R}\left( w\right) $ at a unique point $\bar{w}$ from below.

For \textbf{question b)}, we now have%
\begin{equation*}
w-T\left( w\right) =w-\tau \left( w-a\right) =\left( 1-\tau \right) w+\tau
a,\qquad a>0
\end{equation*}%
and hence%
\begin{equation*}
1-T^{\prime }\left( w\right) =1-\tau .
\end{equation*}%
Observing the curves $h_{L}\left( w\right) $ and $h_{R}\left( w\right) $, an
increase in $a$ shifts $h_{L}\left( w\right) $ upward and leaves $%
h_{R}\left( w\right) $ unchanged. Since $h_{R}\left( w\right) $ is
decreasing and $h_{L}\left( w\right) $ intersects it from below, an upward
shift in $h_{L}\left( w\right) $ shifts the intersect $\bar{w}$ to the left
and the reservation wage decreases. Finally notice that the parameter $a$
does not explicitly appear in equation (\ref{eq:Pissarides_opt1}). This
optimality condition implies a negative relationship between $\bar{w}$ and $%
e^{\ast }$ (in case that nothing else in that equation changes), and so
optimal effort increases.

Given that offers now arrive more frequently ($\pi \left( e^{\ast }\right) $
increases), and each arrived offer is more likely to be accepted\ ($\bar{w}$
decreases), the worker is unconditionally more likely to accept employment.

Observe that $a$ represents an untaxed deduction from wage income (in fact,
an employed worker receives a tax rebate on income $a$), so increasing this
deduction increases incentives to search for and accept employment.

For \textbf{question c)}, assume for simplicity that $a=0$, and denote $\tau
=1-t$.\ The two optimality conditions now read%
\begin{eqnarray*}
1 &=&\frac{\pi ^{\prime }\left( e^{\ast }\right) \left( 1-t\right) }{1-\beta
\left( 1-\alpha \right) }\int_{\bar{w}}^{B}\left( w^{\prime }-\bar{w}\right)
dF\left( w^{\prime }\right)  \\
\left( 1-t\right) \bar{w}+ta-\left( z-\beta e^{\ast }\right)  &=&\frac{\beta
\pi \left( e^{\ast }\right) }{\pi ^{\prime }\left( e^{\ast }\right) }.
\end{eqnarray*}%
From the second equation, we express%
\begin{equation*}
\bar{w}\left( e^{\ast },t\right) =\frac{1}{1-t}\left[ z-\beta e^{\ast }+%
\frac{\beta \pi \left( e^{\ast }\right) }{\pi ^{\prime }\left( e^{\ast
}\right) }-ta\right] 
\end{equation*}%
and treat the first equation with $\bar{w}$ as an implicit function of $%
\left( e^{\ast },t\right) $%
\begin{equation}
1-\beta \left( 1-\alpha \right) =\pi ^{\prime }\left( e^{\ast }\right)
\left( 1-t\right) \int_{\bar{w}\left( e^{\ast },\tau \right) }^{B}\left(
w^{\prime }-\bar{w}\left( e^{\ast },t\right) \right) dF\left( w^{\prime
}\right) \doteq f\left( e^{\ast },t\right)   \label{eq:Pissarides_IFT}
\end{equation}%
We have%
\begin{eqnarray*}
\frac{\partial \bar{w}\left( e^{\ast },t\right) }{\partial e^{\ast }} &=&-%
\frac{\beta }{1-t}\frac{\pi \left( e^{\ast }\right) \pi ^{\prime \prime
}\left( e^{\ast }\right) }{\left( \pi ^{\prime }\left( e^{\ast }\right)
\right) ^{2}} \\
\frac{\partial \bar{w}\left( e^{\ast },t\right) }{\partial t} &=&\frac{1}{%
\left( 1-t\right) ^{2}}\left[ z-\beta e^{\ast }+\frac{\beta \pi \left(
e^{\ast }\right) }{\pi ^{\prime }\left( e^{\ast }\right) }-a\right]  \\
\frac{\partial }{\partial \bar{w}}\int_{\bar{w}}^{B}\left( w^{\prime }-\bar{w%
}\right) dF\left( w^{\prime }\right)  &=&-\left( 1-F\left( \bar{w}\right)
\right) 
\end{eqnarray*}%
and we can apply the implicit function theorem to (\ref{eq:Pissarides_IFT})
to obtain%
\begin{equation*}
\frac{de^{\ast }}{dt}=-\frac{\partial f/\partial t}{\partial f/\partial
e^{\ast }}
\end{equation*}%
where%
\begin{eqnarray*}
\frac{\partial f}{\partial t} &=&-\pi ^{\prime }\left( e^{\ast }\right)
\int_{\bar{w}\left( e^{\ast },\tau \right) }^{B}\left( w^{\prime }-\bar{w}%
\left( e^{\ast },t\right) \right) dF\left( w^{\prime }\right)  \\
&&-\pi ^{\prime }\left( e^{\ast }\right) \left( 1-t\right) \left( 1-F\left( 
\bar{w}\right) \right) \frac{1}{\left( 1-t\right) ^{2}}\left[ z-\beta
e^{\ast }+\frac{\beta \pi \left( e^{\ast }\right) }{\pi ^{\prime }\left(
e^{\ast }\right) }-a\right]  \\
\frac{\partial f}{\partial e^{\ast }} &=&\pi ^{\prime \prime }\left( e^{\ast
}\right) \left( 1-t\right) \int_{\bar{w}\left( e^{\ast },\tau \right)
}^{B}\left( w^{\prime }-\bar{w}\left( e^{\ast },t\right) \right) dF\left(
w^{\prime }\right)  \\
&&+\pi ^{\prime }\left( e^{\ast }\right) \left( 1-t\right) \left( 1-F\left( 
\bar{w}\right) \right) \frac{\beta }{1-t}\frac{\pi \left( e^{\ast }\right)
\pi ^{\prime \prime }\left( e^{\ast }\right) }{\left( \pi ^{\prime }\left(
e^{\ast }\right) \right) ^{2}}
\end{eqnarray*}%
Clearly, $\partial f/\partial e^{\ast }<0$, and a sufficient condition for $%
\partial f/\partial t<0$ is%
\begin{equation*}
z-\beta e^{\ast }+\frac{\beta \pi \left( e^{\ast }\right) }{\pi ^{\prime
}\left( e^{\ast }\right) }-a>0
\end{equation*}%
which is plausible, especially for small $a$. In that case, we indeed obtain
that%
\begin{equation*}
\frac{de^{\ast }}{dt}<0
\end{equation*}%
and an increase in the tax rate reduces effort. Consequently%
\begin{eqnarray*}
\frac{dw\left( e^{\ast },t\right) }{dt} &=&\frac{\partial \bar{w}\left(
e^{\ast },t\right) }{\partial e^{\ast }}\frac{de^{\ast }}{dt}+\frac{\partial 
\bar{w}\left( e^{\ast },t\right) }{\partial t} \\
&=&-\frac{\beta }{1-t}\frac{\pi \left( e^{\ast }\right) \pi ^{\prime \prime
}\left( e^{\ast }\right) }{\left( \pi ^{\prime }\left( e^{\ast }\right)
\right) ^{2}}\frac{de^{\ast }}{dt}+\frac{1}{\left( 1-t\right) ^{2}}\left[
z-\beta e^{\ast }+\frac{\beta \pi \left( e^{\ast }\right) }{\pi ^{\prime
}\left( e^{\ast }\right) }-a\right] .
\end{eqnarray*}%
Substituting in $de^{\ast }/dt$ from the above calculations and simplyfing,
we obtain%
\begin{equation*}
\frac{dw\left( e^{\ast },t\right) }{dt}=\frac{1}{1-t}\frac{\left[ z-\beta
e^{\ast }-a\right] \int_{\bar{w}\left( e^{\ast },\tau \right) }^{B}\left(
w^{\prime }-\bar{w}\left( e^{\ast },t\right) \right) dF\left( w^{\prime
}\right) }{\left( 1-t\right) \int_{\bar{w}\left( e^{\ast },\tau \right)
}^{B}\left( w^{\prime }-\bar{w}\left( e^{\ast },t\right) \right) dF\left(
w^{\prime }\right) +\left( 1-F\left( \bar{w}\right) \right) \frac{\beta \pi
\left( e^{\ast }\right) }{\pi ^{\prime }\left( e^{\ast }\right) }}
\end{equation*}%
which will be positive if and only if%
\begin{equation*}
z-\beta e^{\ast }-a>0\text{.}
\end{equation*}%
Notice that if this condition is satisfied, then also the sufficient
condition for $de^{\ast }/dt<0$ will be satisfied.

Intuitively, an increase in the tax rate will decrease search effort and
increase reservation wages if $a$ is not too high. When $a$ is high, then an
increase in the tax rate $\tau $ substantially increases the rebate $\tau a$%
, and effectively raises post-tax wage income, so that the force operates in
the opposite direction, making employment more attractive.
\end{solution}
\end{exercise}

\begin{draftnote}
\textbf{Calculations for the above problem}: We need to substitute $de^{\ast
}/dt$:%
\begin{eqnarray*}
\frac{dw\left( e^{\ast },t\right) }{dt} &=&\frac{\beta }{1-t}\frac{\pi
\left( e^{\ast }\right) \pi ^{\prime \prime }\left( e^{\ast }\right) }{%
\left( \pi ^{\prime }\left( e^{\ast }\right) \right) ^{2}}\frac{-\pi
^{\prime }\left( e^{\ast }\right) \int_{\bar{w}\left( e^{\ast },\tau \right)
}^{B}\left( w^{\prime }-\bar{w}\left( e^{\ast },t\right) \right) dF\left(
w^{\prime }\right) -\pi ^{\prime }\left( e^{\ast }\right) \left( 1-t\right)
\left( 1-F\left( \bar{w}\right) \right) \frac{1}{\left( 1-t\right) ^{2}}%
\left[ z-\beta e^{\ast }+\frac{\beta \pi \left( e^{\ast }\right) }{\pi
^{\prime }\left( e^{\ast }\right) }-a\right] }{\pi ^{\prime \prime }\left(
e^{\ast }\right) \left( 1-t\right) \int_{\bar{w}\left( e^{\ast },\tau
\right) }^{B}\left( w^{\prime }-\bar{w}\left( e^{\ast },t\right) \right)
dF\left( w^{\prime }\right) +\pi ^{\prime }\left( e^{\ast }\right) \left(
1-t\right) \left( 1-F\left( \bar{w}\right) \right) \frac{\beta }{1-t}\frac{%
\pi \left( e^{\ast }\right) \pi ^{\prime \prime }\left( e^{\ast }\right) }{%
\left( \pi ^{\prime }\left( e^{\ast }\right) \right) ^{2}}} \\
&&+\frac{1}{\left( 1-t\right) ^{2}}\left[ z-\beta e^{\ast }+\frac{\beta \pi
\left( e^{\ast }\right) }{\pi ^{\prime }\left( e^{\ast }\right) }-a\right] 
\end{eqnarray*}%
and simplifying%
\begin{eqnarray*}
\frac{dw\left( e^{\ast },t\right) }{dt} &=&-\frac{1}{\left( 1-t\right) ^{2}}%
\frac{\beta \pi \left( e^{\ast }\right) }{\pi ^{\prime }\left( e^{\ast
}\right) }\frac{\left( 1-t\right) \int_{\bar{w}\left( e^{\ast },\tau \right)
}^{B}\left( w^{\prime }-\bar{w}\left( e^{\ast },t\right) \right) dF\left(
w^{\prime }\right) +\left( 1-F\left( \bar{w}\right) \right) \left[ z-\beta
e^{\ast }+\frac{\beta \pi \left( e^{\ast }\right) }{\pi ^{\prime }\left(
e^{\ast }\right) }-a\right] }{\left( 1-t\right) \int_{\bar{w}\left( e^{\ast
},\tau \right) }^{B}\left( w^{\prime }-\bar{w}\left( e^{\ast },t\right)
\right) dF\left( w^{\prime }\right) +\left( 1-F\left( \bar{w}\right) \right) 
\frac{\beta \pi \left( e^{\ast }\right) }{\pi ^{\prime }\left( e^{\ast
}\right) }} \\
&&+\frac{1}{\left( 1-t\right) ^{2}}\left[ z-\beta e^{\ast }+\frac{\beta \pi
\left( e^{\ast }\right) }{\pi ^{\prime }\left( e^{\ast }\right) }-a\right] 
\end{eqnarray*}%
and again%
\begin{equation*}
\frac{dw\left( e^{\ast },t\right) }{dt}=\frac{1}{1-t}\frac{\left[ z-\beta
e^{\ast }-a\right] \int_{\bar{w}\left( e^{\ast },\tau \right) }^{B}\left(
w^{\prime }-\bar{w}\left( e^{\ast },t\right) \right) dF\left( w^{\prime
}\right) }{\left( 1-t\right) \int_{\bar{w}\left( e^{\ast },\tau \right)
}^{B}\left( w^{\prime }-\bar{w}\left( e^{\ast },t\right) \right) dF\left(
w^{\prime }\right) +\left( 1-F\left( \bar{w}\right) \right) \frac{\beta \pi
\left( e^{\ast }\right) }{\pi ^{\prime }\left( e^{\ast }\right) }}
\end{equation*}
\end{draftnote}

\begin{draftnote}
why not use $c$ for unemployment compensation?
\end{draftnote}

\end{document}
